% ======================================================================
% Plantilla de Informe Técnico / Documento Académico - UTN FRA / LSE
% ======================================================================
\documentclass[DOC]{UTNLSE}

% Rutas para búsqueda de figuras y gráficos
\graphicspath{{images/}{./}}

% ----------------------------------------------------------------------
% METADATOS DEL DOCUMENTO (Opción [RT])
% ----------------------------------------------------------------------
% \titleES{Título del Informe o Trabajo Práctico}
% \subtitleES{Subtítulo o Descripción Breve} % Opcional
% \titleEN{English Title of the Report or Technical Paper}
% \subtitleEN{Optional English Subtitle}

% % Autores: \addauthor{Nombre}{Apellido}{Institución / Cátedra}{ORCID o Enlace}[email]
% \addauthor{Nombre}{Apellido}{Universidad Tecnológica Nacional - FRA}{https://orcid.org/0000-0000-0000-0000}[autor@frba.utn.edu.ar]

% % Resumen / Abstract y Palabras Clave
% \resumenES{Breve resumen del trabajo (100 a 200 palabras). Incluir objetivos, metodología aplicada, resultados destacados y conclusiones principales.}
% \keywordsES{palabra clave 1; palabra clave 2; palabra clave 3}

% \resumenEN{Brief summary of the document in English covering objectives, methodology, main findings, and conclusions.}
% \keywordsEN{keyword 1; keyword 2; keyword 3}

% ----------------------------------------------------------------------
% METADATOS DEL DOCUMENTO (Opción [DOC] - Descomentar si se usa \documentclass[DOC]{UTNLSE})
% ----------------------------------------------------------------------
\title{ROS 2 - Sistema Operativo de Robótica}
\subtitle{Proyecto de investigación - Sistemas de Comunicaciones}
\addauthor{Pablo Sebastian}{Gomez}
\addauthor{Irina}{Terebiznik}
\addauthor{Matias}{Peña}
\course{Sistemas de Comunicaciones}
\coverimage{logo-universidad2.png}
% \LSE

% ======================================================================
% CUERPO DEL DOCUMENTO
% ======================================================================
\begin{document}

\section*{Extracto}
El presente informe analiza ROS (\textit{Robot Operating System}) desde la perspectiva de los sistemas de comunicaciones. Se describen sus conceptos arquitectónicos y patrones de comunicación, y se examina la transición desde la arquitectura híbrida de ROS~1, dependiente de un proceso maestro central, hacia el modelo descentralizado de ROS~2, sustentado en el estándar DDS y en el protocolo RTPS. Se analizan, asimismo, los mecanismos de calidad de servicio (QoS), la capa de abstracción de \textit{middleware} (RMW) y sus principales implementaciones, el soporte para sistemas de tiempo real y microcontroladores, y el esquema de distribuciones del proyecto. El estudio se complementa con una prueba de concepto sobre hardware dedicado.

\section{Capítulo 1}
\subsection{Introducción del trabajo}

El presente trabajo de investigación se desarrolla en el marco de la asignatura
Sistemas de Comunicaciones de la Universidad Tecnológica Nacional. Su propósito
es analizar ROS (\textit{Robot Operating System}) desde la perspectiva de los
sistemas de comunicaciones, considerando los mecanismos de intercambio de
información que emplea y los estándares sobre los que se sustenta.

Los sistemas robóticos actuales constituyen, en esencia, sistemas distribuidos.
Sensores, actuadores, unidades de procesamiento y estaciones de supervisión
operan como procesos independientes, frecuentemente alojados en distintos
equipos, que deben intercambiar datos de manera continua a través de una red.
En este contexto, la arquitectura de comunicaciones deja de ser un aspecto
secundario y pasa a condicionar directamente la latencia, la confiabilidad y la
escalabilidad del sistema en su conjunto.

ROS se ha consolidado como el \textit{framework} de referencia para el
desarrollo de software robótico, tanto en el ámbito académico como en el
industrial. Su evolución resulta particularmente relevante para el estudio de
las comunicaciones: la primera versión (ROS~1) se basa en una arquitectura
híbrida que depende de un nodo maestro central, mientras que ROS~2 adopta un
esquema descentralizado sustentado en el estándar DDS (\textit{Data
Distribution Service}), el protocolo RTPS (\textit{Real-Time Publish-Subscribe})
y el transporte sobre UDP/IP. Este trabajo aborda dicha transición, analizando
sus fundamentos técnicos, la configuración de la calidad de servicio (QoS) y
las implementaciones de \textit{middleware} disponibles.

\subsection{Objetivo y alcance}

El objetivo general del trabajo es analizar el estado del arte de ROS, con foco
en la evolución de su arquitectura de comunicaciones desde ROS~1 hasta ROS~2 y
en los estándares y protocolos que la sustentan.

Para ello se plantean los siguientes objetivos específicos:
\begin{itemize}
    \item Describir los conceptos arquitectónicos de ROS (nodos y \textit{topics})
    y sus patrones de comunicación (\textit{services} y \textit{actions}).
    \item Analizar la arquitectura híbrida de ROS~1 y las limitaciones derivadas
    de su dependencia del nodo maestro (\textit{roscore}) como punto único de
    falla.
    \item Describir el modelo descentralizado de ROS~2 y la pila de
    comunicaciones DDS, RTPS y UDP/IP.
    \item Analizar los mecanismos de calidad de servicio (QoS) y su
    configuración.
    \item Relevar el soporte de ROS~2 para sistemas operativos de tiempo real
    (RTOS) y microcontroladores mediante micro-ROS.
    \item Describir la capa de abstracción de \textit{middleware} (RMW) y
    comparar sus principales implementaciones: Eclipse Cyclone DDS, eProsima
    Fast DDS y RTI Connext DDS.
    \item Presentar el esquema de distribuciones de ROS y su vinculación con
    las versiones de Ubuntu.
    \item Realizar una demostración práctica sobre hardware dedicado (Raspberry
    Pi) que ilustre el intercambio de mensajes entre nodos de ROS~2 a través de
    la red.
\end{itemize}

En cuanto al alcance, el estudio se centra en la capa de \textit{middleware} y
transporte de datos, complementado con una prueba de concepto de carácter
demostrativo. No se abordan en profundidad los algoritmos de
percepción, planificación o control, ni el desarrollo de aplicaciones
robóticas específicas.

\subsection{Metodología y Fuentes de investigación}
La investigación si el trabajo de investigación es del tipo descriptivo y posee
atributos cualitativos o tiene otro foco.

\subsection{Delimitación del estudio}
La delimitación del análisis se da en las áreas por ejemplo de:

- Sector Telecomunicaciones
- Producto: servicio de telefonía móvil

\subsection{Universo Espacio Temporal}
Los años en los cuales se tuvo presente la información para hacer el trabajo.

\section{Capítulo 2}
\subsection{Introducción del tema}

ROS (\textit{Robot Operating System}) es un \textit{framework} de código
abierto para el desarrollo de software robótico. A pesar de su nombre, no
constituye un sistema operativo en sentido estricto: se ejecuta sobre un
sistema operativo anfitrión, principalmente Linux, y provee una capa de
\textit{middleware} junto con un conjunto de bibliotecas y herramientas. Su
función central es permitir que los distintos componentes de un robot se
desarrollen como procesos independientes y se comuniquen entre sí mediante
interfaces estandarizadas.

El proyecto se originó en 2007 en el Laboratorio de Inteligencia Artificial de
la Universidad de Stanford y continuó su desarrollo en Willow Garage.
Actualmente es mantenido por la Open Source Robotics Foundation (OSRF) junto a
una amplia comunidad internacional.

Las principales razones de su adopción son:
\begin{itemize}
    \item \textbf{Modularidad:} cada función del robot se implementa como un
    componente independiente, lo que facilita el desarrollo en paralelo y el
    mantenimiento.
    \item \textbf{Reutilización:} existe un amplio repositorio de paquetes para
    navegación, percepción, manipulación y controladores de hardware.
    \item \textbf{Abstracción de hardware:} los componentes se comunican a
    través de interfaces definidas, independientemente del fabricante del
    sensor o actuador.
    \item \textbf{Herramientas:} incluye utilidades de visualización, registro,
    simulación y depuración.
    \item \textbf{Licenciamiento abierto:} se distribuye bajo licencias
    permisivas (BSD y Apache~2.0), lo que habilita su uso académico e
    industrial.
\end{itemize}

ROS~1 fue concebido como una plataforma de investigación, orientada a un único
robot operando en una red local confiable. A medida que su uso se extendió a
entornos industriales, sistemas multi-robot y redes inalámbricas, sus
limitaciones de diseño se hicieron evidentes. En respuesta, se desarrolló
ROS~2, rediseñado desde su base sobre estándares industriales de
comunicaciones.

La diferencia fundamental entre ambas versiones reside en su arquitectura de
comunicaciones. ROS~1 emplea un esquema híbrido en el que un nodo maestro
central (\textit{roscore}) gestiona el descubrimiento de los participantes,
mientras que el intercambio de datos se realiza mediante protocolos propios
(TCPROS y UDPROS). Este diseño introduce un punto único de falla y no ofrece
mecanismos para ajustar el comportamiento de la comunicación según las
condiciones de la red.

ROS~2, en cambio, adopta una arquitectura descentralizada basada en el
estándar DDS (\textit{Data Distribution Service}) y su protocolo de
interoperabilidad RTPS. Esta elección elimina la dependencia de un nodo
central y habilita la configuración de políticas de calidad de servicio (QoS),
el soporte para sistemas operativos de tiempo real y la integración de
microcontroladores mediante micro-ROS. Asimismo, incorpora mecanismos de
seguridad basados en la especificación DDS-Security, ausentes en ROS~1.

Con la finalización del soporte de ROS~1 en 2025, ROS~2 se ha consolidado como
la plataforma vigente para el desarrollo de nuevos sistemas robóticos.

Los aspectos técnicos de cada arquitectura se desarrollan en detalle en el
Capítulo~3.

\subsection{Descripción del contexto del tema}

ROS se utiliza como capa de integración en sistemas robóticos de diversa
complejidad, desde plataformas educativas y de investigación hasta vehículos
autónomos y robots industriales. En estos sistemas, los componentes de
software suelen distribuirse entre varios equipos: computadoras embarcadas en
el robot, microcontroladores asociados a sensores y actuadores, y estaciones de
supervisión remotas. ROS actúa como el \textit{middleware} que permite que
todos estos componentes intercambien información de forma transparente, sin
importar en qué equipo se ejecuten.

El software de ROS se publica en \textbf{distribuciones}: conjuntos versionados
de paquetes que se prueban y se mantienen en conjunto. Desde 2017, ROS~2 lanza
una nueva distribución cada año, el 23 de mayo. Las distribuciones de los años
pares son de soporte extendido (LTS, \textit{Long Term Support}) y reciben
actualizaciones durante cinco años, mientras que las de los años impares tienen
un ciclo de soporte de aproximadamente un año y medio. En paralelo existe
Rolling, una distribución de desarrollo continuo en la que se incorporan los
cambios antes de su inclusión en una versión estable.

Cada distribución está vinculada a una versión LTS específica de Ubuntu, que
constituye su plataforma principal de soporte (\textit{Tier~1}), según lo
establecido en la REP~2000. Este vínculo se debe a que los paquetes binarios de
ROS se compilan y se prueban contra las bibliotecas del sistema de esa versión
de Ubuntu. En consecuencia, la elección de la distribución de ROS determina la
versión del sistema operativo a utilizar. Las distribuciones LTS más recientes
y su correspondencia son:
\begin{itemize}
    \item \textbf{ROS~1 Noetic Ninjemys} (2020): Ubuntu 20.04. Última
    distribución de ROS~1, con fin de soporte en mayo de 2025.
    \item \textbf{ROS~2 Humble Hawksbill} (2022): Ubuntu 22.04.
    \item \textbf{ROS~2 Jazzy Jalisco} (2024): Ubuntu 24.04.
    \item \textbf{ROS~2 Lyrical Luth} (2026): Ubuntu 26.04.
\end{itemize}

Además de Ubuntu, las distribuciones de ROS~2 ofrecen distintos niveles de
soporte para otras plataformas, como Windows y Red Hat Enterprise Linux.

\subsection{Mercado y Segmentación del tema}

Como es el mercado de usabilidad y la segmentación del mismo. Como se usa
en el mundo y argentina

\subsection{Legislación del tema}
Si existen leyes o normas sobre el tema, el contexto mundial o de argentina

\section{Capítulo 3}
\subsection{Fundamentos técnicos}

\subsubsection{Conceptos arquitectónicos: nodos y \textit{topics}}

Un sistema basado en ROS se organiza como un grafo de procesos que se
comunican entre sí, denominado grafo de ROS. La unidad básica de este grafo es
el \textbf{nodo}: un participante que realiza una función lógica específica,
como leer un sensor, estimar la posición del robot o controlar un motor. Cada
nodo utiliza una biblioteca cliente para comunicarse con el resto del sistema:
\texttt{roscpp} y \texttt{rospy} en ROS~1, y \texttt{rclcpp} y \texttt{rclpy}
en ROS~2, para los lenguajes C++ y Python respectivamente. Los nodos pueden
comunicarse con otros nodos alojados en el mismo proceso, en otro proceso o en
otro equipo de la red. Las conexiones entre ellos se establecen mediante un
mecanismo de descubrimiento, cuya implementación difiere entre ROS~1 y ROS~2,
según se detalla más adelante.

El mecanismo de comunicación fundamental es el \textbf{\textit{topic}}: un canal
con nombre sobre el que opera un modelo de publicación/suscripción. Los nodos
productores de datos (\textit{publishers}) publican mensajes en un
\textit{topic}, y todos los nodos suscritos a él (\textit{subscribers}) los
reciben. Un mismo \textit{topic} admite cero o más publicadores y cero o más
suscriptores, lo que configura un esquema de comunicación de muchos a muchos.
Este modelo presenta dos propiedades relevantes:
\begin{itemize}
    \item \textbf{Anonimato:} el suscriptor no necesita conocer qué nodo
    originó el dato. Esto permite reemplazar, agregar o quitar nodos sin
    modificar el resto del sistema.
    \item \textbf{Tipado fuerte:} cada \textit{topic} transporta un único tipo
    de mensaje, definido en un archivo \texttt{.msg}. A partir de esta
    definición, las herramientas de ROS generan automáticamente el código para
    los distintos lenguajes soportados.
\end{itemize}

Los \textit{topics} resultan adecuados para flujos de datos continuos, como
las lecturas de sensores o el estado del robot, donde no se espera una
respuesta por parte del receptor.

\subsubsection{Patrones de comunicación: \textit{services} y \textit{actions}}

Además de los \textit{topics}, ROS ofrece dos patrones de comunicación
orientados a la interacción entre un cliente y un servidor.

Un \textbf{\textit{service}} implementa una llamada a procedimiento remoto
(RPC, \textit{Remote Procedure Call}) bajo un esquema de solicitud y
respuesta. Un nodo cliente envía una solicitud, el nodo servidor realiza el
procesamiento correspondiente y devuelve un resultado. Su interfaz se define en
un archivo \texttt{.srv}, compuesto por un mensaje de solicitud y uno de
respuesta. Cada nombre de \textit{service} debe tener un único servidor, aunque
puede ser utilizado por cualquier cantidad de clientes. Dado que el cliente
generalmente queda a la espera del resultado, los \textit{services} están
pensados para operaciones de ejecución breve, como consultar un estado o
modificar una configuración.

Una \textbf{\textit{action}} es una llamada a procedimiento remoto de larga
duración, que incorpora la posibilidad de recibir retroalimentación durante la
ejecución y de cancelar el objetivo antes de su finalización. Un ejemplo típico
es la orden de desplazar el robot hasta una posición determinada: la tarea
puede demorar segundos o minutos, durante los cuales el servidor informa el
progreso y el cliente puede cancelarla. Su interfaz se define en un archivo
\texttt{.action}, compuesto por tres mensajes: objetivo (\textit{goal}),
resultado (\textit{result}) y retroalimentación (\textit{feedback}).

En ROS~2, las \textit{actions} no constituyen un mecanismo de transporte
independiente, sino que se construyen sobre los dos patrones anteriores: cada
\textit{action} se compone internamente de tres \textit{services} (envío del
objetivo, cancelación y obtención del resultado) y dos \textit{topics}
(retroalimentación y estado). En ROS~1, en cambio, las \textit{actions} no
formaban parte del núcleo y se implementaban mediante una biblioteca
independiente denominada \texttt{actionlib}.

\subsection{Ámbitos de aplicación}

Donde se utiliza y como, en que mercados, etc.

\subsection{Estrategias Competitivas}
¿Competencia? ¿Producto único?

\subsection{Normativas asociadas}\label{sec:normativas}
¿Qué normas tiene asociadas ITU, IEEE?

\section{Capítulo 4}
\subsection{Conclusiones}

Al finalizar el ...

\subsection{Glosario de términos}
Terminología utilizada en el informe:
\begin{itemize}
    \item \textbf{ROS 2:} \textit{Robot Operating System 2}. Segunda generación de ROS, \textit{framework} de código abierto para el desarrollo de software robótico.
    \item \textbf{Nodo:} Participante del grafo de ROS que realiza una función lógica específica y se comunica con otros nodos.
\end{itemize}

% ======================================================================
% BIBLIOGRAFÍA EMBEBIDA
% ======================================================================
\subsection{Referencias}
% thebibliography crea su propio \section*{Referencias}; se anula dentro
% del grupo para que quede como subsección 4.3, igual que la plantilla.
\begingroup
\renewcommand{\section}[2]{}
\begin{thebibliography}{99}

\bibitem{ref_ejemplo1}
A. S. Sedra y K. C. Smith, \emph{Microelectronic Circuits}, 7ma ed. Oxford University Press, 2015.

\bibitem{datasheet_ejemplo}
Texas Instruments, \emph{LM317 3-Terminal Adjustable Regulator}, Datasheet SLVS044X, 2020. [Disponible en carpeta \texttt{references/}].

\end{thebibliography}
\endgroup

\end{document}
